% Section 3 fragment for the all-in-one manuscript.
% Outline mapping: P-3.01--P-3.14. Preamble and review wrapper: main.tex.
\section{Shared Objects and Contracts}
\label{sec:shared-objects}

This section fixes the objects exchanged by the protocol and the state
that each component owns. A session's committee and initial anchor are
immutable inputs to its regional engine. Public replay determines which
regional results may extend the global chain; the Session Manager
controls the lifetime of local engine instances. These responsibilities
use the same identifiers and evidence types.

\subsection{Sessions, policies, and contexts}
\label{sec:objects-context}

% P-3.01
\begin{definition}[Session and view]
\label{def:session-view}
A \emph{session} is a pair $s=(\beta,r)$, where $\beta\in\mathbb N$ is a
branch version and $r\in\mathbb N$ is a turn index. Once a region is
selected, it is denoted by $i_s$ and its fixed committee by $V_s$.
A regional view is
$(\beta,r,g,t)$: $g\in\mathbb N$ is a recovery generation and
$t\in\mathbb N$ is a leader term within that generation. The index $j$
names a fresh proposal position within a term.
\end{definition}
A term change advances $t$; a Quit may open $g+1$ and restart at $t=0$.
A regular handoff opens the next turn on the same branch. An executing
Force or RESUME opens a new branch version, while a STOP closes the
current ordinary-production branch without itself opening a regional
session. A new version may reuse a turn index, so $r$ alone never
identifies an engine. Each replica keeps independent mutable state for
each session in which it participates.
Subscripts such as $A_{s,g,t}$ abbreviate $A_{\beta,r,g,t}$ when
$s=(\beta,r)$.

\begin{table}[htbp]
\centering\small
\caption{Principal objects and their scope. A field fixed by a session is
not rewritten by a later snapshot, recovery, or mode change.}
\label{tab:notation}
\begin{tabularx}{\linewidth}{@{}P{.31\linewidth}Y@{}}
\toprule
Notation & Meaning\\\midrule
$s=(\beta,r)$; $(\beta,r,g,t),j$ & Session; regional view and fresh position.\\
$i_s,V_s,w_o,W_s,F_s$ & Region, committee, immutable weights, their total,
and the session's fault bound.\\
$\mathcal C_e$; $a_e,c_e$ & Staking snapshot, activation height, and cutoff
of staking epoch $e$.\\
$\mathcal P_\beta=(S_\beta,\mathcal L_\beta,\lambda_\beta)$ & Immutable
branch policy: region status map, cyclic order, and L1 floor.\\
$\Pi_s$; $\mathsf{AUTH}_{s,g}$ & Session plan; generation authority.\\
$E_s$; $A_s=A_{\beta,r,0,0}$ & Selection reference; initial execution
anchor of the session.\\
$C_s=(\Pi_s,A_s)$; $\mathrm{cid}_s$ & Session context and its identifier.\\
$A_{\beta,r,g,t}$ & Anchor preceding the fresh positions of this view.\\
$h(B)$; $k_s(B)$ & L2-block height; physical regional position, counting
DATA and DRAIN.\\
$\TB,K,H_s,T_s$ & Session DATA quota, term limit, terminal DATA height,
and terminal DATA on the session ancestry.\\
$\QC_{s,g,t}(B)$; $\FC_s(B)$ & Proposal certificate; regional FINAL
receipt or stated same-session covering evidence.\\
$Q_s^0,Q_v,Q^*$ & Inherited base justification, entry lock, and selected
recovery justification.\\
$\rho(B),\lambda(B)$; $\xi(B)$ & Finalized L1 reference and its height;
finalized Espresso header reference.\\
$J_q,\theta_q$; $\tau$ & Espresso block and timestamp; real time.\\
$q_s^{\rm open},\lambda_s^{\rm open}$ & Lower bounds on references in
the session's first DATA.\\
$\nu,\omega,\sigma_G(B)$ & Selected-session descriptor, unassigned
boundary, and immutable global provenance of a block.\\
$G,R,\mathsf{Tail}$ & Global chain, regional production branch, and
responsibility for the next global block.\\
$\mu,\kappa,pc$ & Production mode, control-processing state, and durable
replay execution position.\\
\bottomrule
\end{tabularx}
\end{table}

% P-3.02
\begin{definition}[Branch policy and session plan]
\label{def:policy-plan}
Each branch version has one policy
\[
 \mathcal P_\beta=(S_\beta,\mathcal L_\beta,\lambda_\beta).
\]
Here $S_\beta$ assigns exactly one of $\mathsf{ON}$ or $\mathsf{OFF}$ to
each registered region, $\mathcal L_\beta$ is a cyclic permutation of
those regions, and $\lambda_\beta$ is the branch's L1 floor. The policy
is installed at genesis or by an executed branch-opening control and
remains fixed throughout that version. Floors do not decrease across
executed branch openings.
For a target session $s=(\beta,r)$, the Authority derives
\[
 \Pi_s=\mathsf{Plan}(\mathcal P_\beta,\beta,r,E_s,i_{\rm src})
\]
from its fixed selection reference $E_s$ and source region $i_{\rm src}$,
with $i_{\rm src}=\bot$ at genesis.
A plan records these inputs, the authenticated snapshot and initial ban
set used for selection, the selected region or $\bot$, a source height
$\lambda(\Pi_s)$, and a deterministic identifier $\mathrm{id}(\Pi_s)$.
If a region is selected, it also supplies the committee and weights,
$W_s,F_s,\TB,K$, and $\Delta_{i_s}$. A $\bot$ result supplies no
committee or regional engine.
\end{definition}
The plan binds its selection inputs and results through canonical
encodings and authenticated source identifiers. The source height $\lambda(\Pi_s)$ is the greatest L1 publication or
activation height required by its source records, not the reader's
current L1 head. It does not contain the execution anchor. For a regular successor $s'$ of $s$, selection uses
$E_{s'}=A_s$ while execution starts after $A_{s'}=T_s$. For the first
session of a branch, the selection reference and execution anchor
coincide. Thus a regular successor can be planned before its exact
execution anchor becomes available. The Authority computation fixes the
snapshot and ban cut; missing source records yield \textsf{WAIT}, not a
zero-weight region or a different snapshot. The L1 floor is a lower
bound, rather than necessarily the inclusion height of the most recent
control when controls span a mode transition.

For a selected session, the initial authority is
$\mathsf{AUTH}_{s,0}=\Pi_s$. For $g>0$,
$\mathsf{AUTH}_{s,g}$ augments the plan with the unique finalized recovery
record $\mathsf{REC}(\mathrm{cid}_s,g)$ and its fixed ban set. The
Authority supplies the deterministic leader
\[
 L(\beta,r,g,t)=\mathsf{Leader}(\mathsf{AUTH}_{s,g},t).
\]
Recovery changes proposer eligibility, not the committee, voting weights,
quota, or threshold. A member banned from proposing remains a voting
member with its original session weight. A verified empty proposer set
returns $\bot$; an unavailable authority returns \textsf{WAIT}.

% P-3.03
\begin{definition}[Session context]
\label{def:session-context}
For a plan $\Pi_s$ that selects a region, a session context is
\[
 C_s=(\Pi_s,A_s),\qquad
 \mathrm{cid}_s=\Hash\bigl(\mathrm{id}(\Pi_s),\Hash(A_s)\bigr),
 \qquad A_s=A_{\beta,r,0,0}.
\]
$A_s$ denotes the exact initial anchor, including its original identity;
it is not the opening selected by a later recovery. The plan and anchor
together determine the chain, session, region, fixed committee and
weights, initial execution point, and DATA endpoint.
\end{definition}
An \emph{authentic incoming context} is derived by checking the handoff
history from the authorized starting checkpoint and recomputing its
plan from the required sources. It may belong to a valid speculative
regular successor whose predecessor has not yet entered $G$.
Authenticity therefore establishes historical assignment, not current
permission to run or global acceptance. In contrast, the L1 Authority
checks offense and Quit submissions \emph{relative to a claimed context}:
it checks the fixed sources, signatures, and slots of that context
without deciding whether its session actually opened. A claim cannot
reuse signatures bound to a different $\mathrm{cid}$. Every use that
requires an authentic handoff must still establish that handoff.

\subsection{Blocks, references, and the session endpoint}
\label{sec:objects-blocks}

% P-3.04
\begin{definition}[Regional block and native projection]
\label{def:regional-block}
A fresh regional block has the form
\[
 B=(\beta,r,k,h,\mathit{type},x,\Hash(B_{\rm parent}),\rho,\xi),
 \qquad\mathit{type}\in\{\mathsf{DATA},\mathsf{DRAIN}\}.
\]
For session $s$, set $k_s(A_s)=0$ locally, without modifying the inherited
anchor. Every fresh block satisfies $k_s(B)=k_s(B_{\rm parent})+1$.
A DATA block carries a full application payload $x$, binds one complete
native L2 block, and has $h(B)=h(B_{\rm parent})+1$.
That native block may contain several mini-blocks, whose ordinary
transactions concatenate to its one ordered ordinary-transaction list.
A DRAIN carries no new application input and has
$h(B)=h(B_{\rm parent})$.
\end{definition}
Write $\operatorname{native}(B)$ for the native block bound by a DATA or
L1-derived deposit-only block, and $\mathit{st}(B)$ for the executable
state after that block. DRAIN creates no native block; it preserves the
state of the last DATA on its regional ancestry. The hash $\Hash(B)$ of
the protocol envelope and the native block hash are distinct typed
values. A DATA block binds both its protocol parent and the corresponding
native parent through its payload and native validation. Correct
validation checks both relations; a native hash cannot substitute for
the envelope hash in a protocol signature or parent link.

The global height $h$ counts actual L2 blocks, including deposit-only
blocks in recovery mode. DRAIN advances only $k_s$. A deposit-only block
is derived from finalized L1 inputs, has an empty \emph{ordinary}
transaction list, and can still contain required deposits and system
transactions. Its provenance is an L1 derivation position and mode
boundary, not a regional session or FINAL receipt. Its envelope supplies
an authenticated L1 reference for subsequent planning, but no regional
Espresso reference. It can be an initial anchor of a resumed session;
that new session does not relabel it as its own DATA.

% P-3.05
\paragraph{Finalized references.}
The field $\rho$ identifies a finalized L1 block by height $\lambda(B)$
and hash, together with authenticated access to the required registry
and control prefix. The field $\xi(B)=(q,\Hash(J_q))$ identifies a
finalized Espresso header. Along DATA ancestry, neither referenced
height decreases. The first DATA of $s$ satisfies
\[
 \lambda(B)\ge\lambda_s^{\rm open},\qquad
 \operatorname{height}(\xi(B))\ge q_s^{\rm open}.
\]
A DRAIN copies the L1 and Espresso references of the terminal DATA on
its ancestry. Regional validation checks the referenced facts and
bounds; it applies no receive-time freshness test. The L1 header's
timestamp is not the regional commit time, and the native block's
\texttt{epoch\_num} is an L1-origin height rather than a staking epoch
or an alias for $\lambda(B)$. Global acceptance separately requires the
Espresso reference to be strictly earlier than the Espresso block at
whose replay step the DATA is appended, even when its evidence was
published in earlier blocks. Historical references do not change a session's
fixed voting authority.

% P-3.06
\paragraph{Quota and fresh positions.}
Genesis fixes $\TB\ge1$ and $K\ge1$. For session $s$, the endpoint
\[
 H_s=h(A_s)+\TB
\]
is fixed by its initial anchor. $T_s$ denotes its terminal DATA at that
height on the relevant ancestry; the height is fixed before the block's
identity is known. Recovery does not change $H_s$. Quota counts fresh
DATA in the proposed ancestry, including uncommitted DATA, rather than
only locally committed or globally accepted blocks.

If $h(A_{s,g,0})<H_s$, generation $g$ is a DATA generation. Term $t$
continues after its fixed anchor $A_{s,g,t}$ and has
\begin{equation}
\label{eq:term-positions}
 m_{s,g,t}=\min\{K,H_s-h(A_{s,g,t})\},\quad
 0\le j<m_{s,g,t},
\end{equation}
with positions
\[
 h=h(A_{s,g,t})+j+1,\qquad
 k=k_s(A_{s,g,t})+j+1.
\]
The last fresh position is terminal for the term. Its QC opens the next
term only if its DATA height is strictly below $H_s$; the next term
inherits that block as its anchor. Thus only the last DATA term may
have fewer than $K$ positions, and a QC at $H_s$ opens no later term.

If a recovery generation $g>0$ opens at height $H_s$, it is a
\emph{drain generation}: $t=j=0$, $m_{s,g,0}=1$, and its only fresh block
is a DRAIN justified by the opening's ANCHOR QC. ANCHOR voting starts no
commit timer; the DRAIN vote supplies the timer that can commit the
terminal DATA as an ancestor. A DATA generation does not automatically
append DRAIN after reaching its quota. Further interrupted recoveries
may create further DRAINs, including on a DRAIN opening, but none creates
a new native block or another DATA allowance.

\begin{figure}[tbp]
\centering
\begin{tikzpicture}[x=1cm,y=1cm,font=\small,
 box/.style={draw,rounded corners=2pt,minimum height=.65cm,
             minimum width=3.05cm,align=center,fill=blue!3},
 arr/.style={-{Stealth[length=2mm]},semithick}]
\node[anchor=west,font=\bfseries\small] at (0,3.05)
 {(a) Index hierarchy};
\node[box] (br) at (1.65,2.35) {$\beta$\\branch};
\node[box] (se) at (5.25,2.35) {$(\beta,r)$\\session};
\node[box] (ge) at (8.85,2.35) {$(\beta,r,g)$\\generation};
\node[box] (te) at (12.45,2.35) {$(\beta,r,g,t)$\\term};
\draw[arr] (br)--(se);\draw[arr] (se)--(ge);\draw[arr] (ge)--(te);
\node[align=center,font=\footnotesize] at (5.25,1.60)
 {fixed committee and initial anchor};
\node[align=center,font=\footnotesize] at (12.45,1.60)
 {fresh positions indexed by $j$};
\node[anchor=west,font=\bfseries\small] at (0,1.05)
 {(b) DATA height and physical position: $\TB=10$, $K=4$};
\node[anchor=east,font=\footnotesize] at (1.0,-.05) {$h$:};
\node[anchor=east,font=\footnotesize] at (1.0,-.80) {$k_s$:};
\node[draw,minimum width=.72cm,minimum height=.45cm,fill=black!5]
 (an) at (1.6,-.05) {100};
\node[font=\footnotesize] at (1.6,-.80) {0};
\foreach \i in {1,...,10}{
  \pgfmathsetmacro{\xx}{1.6+.92*\i}
  \pgfmathtruncatemacro{\hh}{100+\i}
  \node[draw,minimum width=.72cm,minimum height=.45cm,fill=teal!8]
   (b\i) at (\xx,-.05) {\hh};
  \node[font=\footnotesize] at (\xx,-.80) {\i};
}
\draw[decorate,decoration={brace,amplitude=3pt},draw=rrBlue]
 (2.16,.28)--(5.64,.28) node[midway,above=4pt,font=\footnotesize] {$t=0$};
\draw[decorate,decoration={brace,amplitude=3pt},draw=rrBlue]
 (5.84,.28)--(9.32,.28) node[midway,above=4pt,font=\footnotesize] {$t=1$};
\draw[decorate,decoration={brace,amplitude=3pt},draw=rrBlue]
 (9.52,.28)--(11.16,.28) node[midway,above=4pt,font=\footnotesize] {$t=2$};
\node[draw,minimum width=1.2cm,minimum height=.45cm,fill=orange!10]
 (dr) at (13.35,-.05) {DRAIN};
\node[font=\footnotesize] at (13.35,-.80) {11};
\node[font=\footnotesize,align=center,text width=2.5cm] at (13.35,-1.35)
 {$h=110$\\no native block};
\draw[arr,draw=rrOrange] (b10.east)--(dr.west)
 node[midway,above=2pt,font=\scriptsize,align=center] {recovery\\at $H_s$};
\node[font=\footnotesize,align=center,text width=10.2cm] at (5.8,-1.35)
 {The initial DATA generation ends at 110. DRAIN belongs to a later recovery.};
\end{tikzpicture}
\caption{Coordinates and quota. Panel (a) shows scopes, not a timing
trace. Panel (b) shows the initial generation's three DATA terms and a
possible later drain generation. The latter advances the physical
position but preserves the DATA height and executable state.}
\label{fig:indices-quota}
\end{figure}

\subsection{Signed messages, certificates, and admission}
\label{sec:objects-messages}

% P-3.07
Every signed regional message authenticates a canonical encoding of its
chain and rule version, message type, $\mathrm{cid}_s$, and typed payload.
The payload contains the session/view coordinates required by its type;
Table~\ref{tab:signed-messages} records the additional fields. All
validation uses the exact producing context. The block tuple in
Definition~\ref{def:regional-block} does not itself contain a plan or
$\mathrm{cid}$: those are bound by the signatures, not inferred from the
block hash alone.

\begin{table}[htbp]
\centering\small
\caption{Authenticated regional messages. Every row also binds the chain,
rule version, type, and session context. An opening identifies the exact
generation-opening block, including its height and physical position.}
\label{tab:signed-messages}
\begin{tabularx}{\linewidth}{@{}P{.21\linewidth}P{.20\linewidth}Y@{}}
\toprule
Type & Coordinates & Payload bound by the message\\\midrule
PROPOSE & $g,t,j$ & Generation opening; full block choice, including type,
$k,h$ and hash; parent justification.\\
DATA-VOTE / DRAIN-VOTE & $g,t,j$ & Generation opening, chosen block type,
$k,h$ and hash. Both types share one fresh-vote slot.\\
STATUS & $g,0$ & Fixed entry lock $Q_v$ and entry proof.\\
NEWVIEW & $g,0$ & Selected opening, justification $Q^*$ and entry proof.\\
ANCHOR-VOTE & $g,0$ & Selected recovery opening.\\
BLAME & $g,t$ & The admitted view; no opening, block hash, or local
progress count.\\
FINAL & $s,h(B)$ & Exact DATA hash. No generation or term field.\\
\bottomrule
\end{tabularx}
\end{table}

\begin{definition}[Regional certificate]
\label{def:regional-certificate}
A regional certificate consists of matching valid signatures from
\emph{distinct} members of $V_s$ whose total session weight is greater
than $F_s$. A proposal QC $\QC_{s,g,t}(B)$ certifies DATA-VOTE or
DRAIN-VOTE messages; an ANCHOR QC certifies the recovery opening; a BLAME
certificate certifies a view complaint; and a FINAL certificate
certifies a DATA hash under its session context. Repeated signatures by
one member contribute its weight only once.
\end{definition}
A correct member uses at most one fresh vote for a slot $(s,g,t,j)$,
including across DATA/DRAIN types, and at most one ANCHOR choice for a
recovery generation. Signatures are kept with their immutable slot
records. A QC establishes certified support, not expiration of a commit
timer. Neither assembling a QC nor casting an ANCHOR vote starts such a
timer. Since FINAL has no $g,t$ field, matching FINAL signatures on one
DATA may combine across generations of the same session, but not across
contexts or sessions.

% P-3.08
\paragraph{Anchors, openings, and justifications.}
The anchor $A_{s,g,t}$ precedes the fresh positions of its view. At
initial entry it is $A_s$; at recovery entry it is the selected opening
$A_{s,g,0}$; and for $t>0$ it is the preceding term's terminal DATA.
The inherited base justification $Q_s^0$ identifies exactly $A_s$ and
its original commitment evidence. If this evidence is a covering
receipt rooted at a later block, that root is evidence only: $Q_s^0$
still justifies $A_s$. For genesis or an L1-derived anchor, the
corresponding checkpoint or derivation evidence has the same base role.
It is not a new vote or QC from the entering committee.

Within one session, a proposal QC for $B$ has lexicographic rank
$(g,t,k_s(B))$. An ANCHOR QC has rank $(g,0,k_s(A_{s,g,0}))$, below the
first fresh QC of that recovery; $Q_s^0$ ranks below all own-session
QCs. A fixed order on certified block/opening identities breaks equal
ranks, independently of signature-set encodings. The opening identity
is $(g,k_s(A_{s,g,0}),h(A_{s,g,0}),\Hash(A_{s,g,0}))$.
$Q_v$ denotes a replica's entry lock, fixed when it enters a recovery
generation, and $Q^*$ the leader's selected justification. Ranks are
compared only within the session; a foreign session's QC cannot replace
$Q_s^0$ or become an own-session entry lock.

\begin{table}[htbp]
\centering\small
\caption{Parent and justification for a fresh proposal. The recovery row
also applies to a drain generation, including one whose opening is an
older DRAIN. No rule requires the parent to be committed before proposing.}
\label{tab:parent-rule}
\begin{tabularx}{\linewidth}{@{}P{.22\linewidth}Y Y@{}}
\toprule
Case & Parent block & Required justification\\\midrule
$g=t=j=0$ & Fixed initial anchor $A_s$ & $Q_s^0$.\\
$g>0$, $t=j=0$ & Recovery opening $A_{s,g,0}$ & Its ANCHOR QC in the
current generation.\\
$j=0$, $t>0$ & Terminal DATA of term $t-1$ & Its proposal QC in that
preceding term.\\
$1\le j<m_{s,g,t}$ & Fresh block at $j-1$ of the same term & Its proposal
QC in the same view.\\
\bottomrule
\end{tabularx}
\end{table}

% P-3.09
\begin{definition}[Admission and complete Quit evidence]
\label{def:admission-quit}
Initial view $(s,0,0)$ is admitted by guarded Enter using its fixed
handoff. A recovery view $(s,g,0)$, $g>0$, requires the preceding
complete Quit proof and finalized $\mathsf{REC}(\mathrm{cid}_s,g)$.
A later term additionally requires the terminal QCs of its preceding
terms, each for DATA below $H_s$. For remote verification, the
\emph{admission witness} supplies the corresponding boundary, recovery
entry proofs and authority records, and terminal-QC ancestry.
A view message is \emph{admissible} when its view is admitted, its
signed context and generation match, and its required signatures are
valid under that context. Leader-signed messages must be signed by
$L(\beta,r,g,t)$ for that exact view.
A complete Quit proof for an admitted generation is one of:
\begin{enumerate}[label=(\alph*)]
\item a BLAME certificate for an admitted term, with that term's
admission witness;
\item two different admissible DATA or DRAIN proposals signed by the
designated leader at the same $(s,g,t,j)$, with the material needed to
check their admission and quota; or
\item for $g>0$, two valid NEWVIEW messages from its designated recovery
leader selecting different opening blocks.
\end{enumerate}
\end{definition}
Different openings conflict even when one extends the other; different
QC encodings for the \emph{same} opening do not constitute that conflict.
An over-quota DATA or a DRAIN outside a drain generation is not an
admissible proposal for case (b). A silent leader does not postpone
admission of its term or the start of its progress checks.

Complete regional validation checks the signatures and authority,
admission, type, quota, both positions, references, parent rule and
justifications, hash links, and native input validity over the required
ancestry. It uses actual block bodies and verified source records,
rather than only their names or hashes. Regional replicas perform this
validation before their own protocol actions. The Authority's relative
checks use the claimed context's fixed sources and signed slots; they
do not establish canonical handoff history or execute the global rule.
The public FINAL kernel is a separate check whose soundness must be
proved from actual correct signing and commit events, not assumed as
part of raw signature verification.

\subsection{Handoffs, readiness, and regional results}
\label{sec:objects-handoffs}

% P-3.10
\begin{definition}[Handoff descriptor]
\label{def:handoff}
A selected session has descriptor
\[
 \nu=(\beta,r,i,A,H,q^{\rm open},\mathit{cause},\Pi),
 \qquad i=\mathsf{region}(\Pi),\quad H=h(A)+\TB.
\]
Its cause is GENESIS, REGULAR, FORCE$(c,q_f)$, or
RESUME$(c,p_{\rm eff})$. Here $c$ is the unique L1 control sequence
number, $q_f$ is the ordinary Force's Espresso activation height, and
$p_{\rm eff}=(\ell,txIndex,logIndex)$ is the finalized L1 position of
the corresponding ModeApplied event. A complete handoff additionally supplies the
branch policy, the selection reference and source region used by the
plan, the L1 bound
\[
 \lambda_s^{\rm open}=\max\{\lambda(A_s),\lambda(\Pi_s)\},
\]
and the evidence fixing these fields and $Q_s^0$.
An unassigned boundary is
\[
 \omega=(\mathit{cause},\beta,r,i_{\rm src},A,q^{\rm open});
\]
it has a fixed target pair and anchor but no selected region.
\end{definition}
A regular handoff fixes $A_{s'}=T_s$, $E_{s'}=A_s$, and
\begin{equation}
\label{eq:regular-reference-floor}
 q_{s'}^{\rm open}=
 \max\{\operatorname{height}(\xi(T_s)),q_s^{\rm open}+1\}.
\end{equation}
Its evidence includes the authentic predecessor boundary and complete
terminal result. A Force handoff uses the exact retained global tip
and its activated control; a resumed handoff uses the actual derived
terminal state of the deposit-only interval and its committed mode
boundary. Their public replay fixes the target pair, source region and
reference floors. Neither the receiver's current local tip nor a
sender's claimed highest production determines these fields. The
inherited anchor retains its original producer or L1 provenance.

Historical handoff verification and current entry permission are
separate predicates. Historical verification reconstructs the fixed
fields from an authorized checkpoint and the complete relevant source
history. Current permission additionally requires the assigned member,
a current version not marked closed in the node's verified history,
and a pair that has not already entered or terminated. The same test
applies if the first received ReadyBundle would create the local
instance. A later source record cannot replace a previously fixed
handoff for the same pair with different fields.

% P-3.11
\begin{definition}[Ready]
\label{def:ready}
An Assigned instance for $s$ is \emph{Ready} when:
\begin{enumerate}[label=(\roman*)]
\item its handoff has been authenticated and current entry permission
holds;
\item the exact anchor, its original identity, required ancestry and
commitment or derivation evidence have been validated;
\item the node holds a validated, directly usable executable state
$\mathit{st}(A_s)$, with no remaining reconstruction or installation
work; and
\item its plan, policy, quota and reference floors, with their required
source records, are available and validated.
\end{enumerate}
\end{definition}
A ReadyBundle carries that handoff, full anchor/evidence material,
executable state and its native validation evidence, and the required
source records. It carries no inherited locks, votes or timers.
Recipients verify its contents and may reuse local state only when it
is the state after the same exact anchor. A state root, URL, or promise
to fetch later is insufficient. The header at $q_s^{\rm open}$ need not
yet exist for Ready or initial Enter; the first DATA must wait until
its own reference condition is satisfiable. ReadyBundle is a delivery
object, not a READY quorum certificate, and a duplicate does not reset
entry time.

% P-3.12
\begin{definition}[Regional FINAL receipt and complete result]
\label{def:regional-result}
Let $C_s$ be an authentic producing context. A direct regional FINAL
receipt $\FC_s(B)$ consists of matching FINAL signatures on the exact
DATA block $B$ of $s$ from distinct members of weight greater than $F_s$;
$B$ extends
$A_s$, lies within its session's quota, and has intrinsically valid
regional ancestry under that context. A covering receipt for $B$ is a
valid direct receipt on a same-session DATA descendant $D$, together
with the exact validated ancestry from $D$ to $B$.
A complete regional result contains the target DATA, the direct or
covering receipt with its full bodies, and the historical context and
source material needed for the consumer's specified validation.
\end{definition}
The reusable public kernel for a target is $(D,\Sigma,\pi)$: root DATA,
FINAL signatures on that root, and the full hash-linked path to the
target. Direct evidence takes $D=B$. Its raw threshold test checks only
the exact signed context, DATA identity and weight; it does not assume
the intrinsic validity claimed in Definition~\ref{def:regional-result}.
The later public-check argument establishes that validity through a
correct signer. Regional full validation additionally uses the actual
proposals, QCs, admission witnesses and finalized Authority records
that its own actions require; these objects do not become mandatory
sidecars of the public kernel merely because a publisher attaches them.

A correct replica emits a new FINAL only when a commit-timer event
commits the DATA directly or as an ancestor. A receipt-induced commit,
including Completion, emits no new FINAL. A DRAIN timer may cause FINAL
for the terminal DATA, never for the DRAIN itself. A successor never
re-signs its inherited anchor under its new committee. Coverage never
crosses sessions: a successor's receipt cannot replace the predecessor's
terminal evidence. A complete receipt may support a regular handoff
before it supports any global append. Regional finality is therefore a
statement about a specified, authenticated production ancestry; it is
not a claim that the ancestry has been selected into $G$.

\subsection{State ownership and invocation rules}
\label{sec:objects-state}

% P-3.13
\begin{table}[htbp]
\centering\small
\caption{State ownership. Immutable validated objects may be shared;
mutable consensus state belongs to exactly one local session instance.}
\label{tab:state-ownership}
\begin{tabularx}{\linewidth}{@{}P{.21\linewidth}Y@{}}
\toprule
Owner & Authoritative local or ledger state\\\midrule
L1 Authority & Positions and escrow, frozen epoch tables, ban history,
REC, processed offenses, control log and clearance/release records.\\
Session Program & Fixed context and endpoint; current view, individual
signing slots, locks, votes, QCs, timers, commitments and execution state.\\
Session Manager & Handoffs, entry/termination records, Closed versions,
production branch $R$, replay state for $G$, and exact prepared anchors.\\
Publisher & Complete-result store, deduplication identifiers, durable
pending sends and publication progress.\\
Batcher/inbox & Local channel construction and, separately, L1-accepted
frame/block progress and committed mode boundaries.\\
Trusted certifier & Validated source/replay state and its typed business
conclusions; no authority to invent a different protocol interpretation.\\
Derivation node & Actual native head and state, consumed L1 positions,
accepted channel data, mode state and pending input dependencies.\\
\bottomrule
\end{tabularx}
\end{table}
The Session Manager follows the global chain $G$ and a possibly longer
regional production branch $R$. It keeps a baseline state at
$\operatorname{last}(G)$ and may prepare a regular successor at a later
regionally finalized anchor. Each Program receives the state of its
own exact anchor. Subsequent baseline progress does not change that
Program's starting point. A shared PublicStore retains authenticated
source bytes, regional evidence, checkpoints and outstanding delivery
tasks independently of Program lifetime. Closing an instance removes
its permission to act, not the obligation to retain or transfer data.

For public replay, $\mathsf{NEXT}(\nu)$ identifies a selected next-block
owner; $\mathsf{PENDING}(\omega)$ records a boundary whose complete plan
selects no region. $\mathsf{OPEN}(\omega,E)$ records fixed inputs whose
plan is not yet available, and $\mathsf{DEPOSIT}$ carries the committed
L1-mode continuation state. OPEN is a transient evidence wait, not a
third production mode. The two modes remain NORMAL and DEPOSIT\_ONLY.
The global rule and derivation rule are deterministic computations
performed by their consumers, not additional consensus services.

% P-3.14
\paragraph{Results, errors, and durable effects.}
A successful query returns a value. \textsf{WAIT} identifies unavailable
required information and leaves the caller at its saved operation;
\textsf{REJECT} rejects a decidably invalid request without its proposed
state change. A verified $\bot$ or empty set is a result obtained from
complete evidence, not a substitute for a failed fetch.
\textsf{INCONSISTENT} is reserved for contradictory authenticated
canonical inputs or a violated continuation invariant; unverified
network garbage cannot trigger it merely by asserting a competing
boundary. A paused caller preserves already committed output and
resumes from its durable execution position.

An authenticated L1 call binds its caller, stable request identifier,
chain, version and arguments; repeated identifiers with identical
content return the recorded result, while conflicting reuse is rejected.
Consensus and publication objects have their own stable content and
slot identities. Changing a QC's signer set or proof encoding does not
create a new block choice or restart a vote timer. State changes, signed
slot records, execution positions and pending outgoing actions are
persisted before dispatch. Retransmission may repeat bytes but not the
logical action. Inputs contain the exact required bodies and source
records, or references that the receiver can resolve within the claimed
processing budget. These conventions apply to every later protocol
block and to the detailed internal interfaces.
